\section{Confidentiality Travels Along the Path, Not the Box}

\textbf{Why this matters:} Sensitive information doesn't stay contained to the step that first touched it --- it travels through logs, caches, and indexes that a firm's existing permission structure was never built to reach.

\textbf{Practical implication:} Ask not whether a vendor isolates your data from other customers, but whether your own internal matter barriers are preserved everywhere the data and its derivatives actually go.

\textbf{The core realization: sensitivity is a property of the path data takes through a system, not a property of any single step.} Once sensitive data enters context assembly, everything downstream inherits its sensitivity --- the model reading it, any memory holding it, any hidden model grading the output, the logs, the step that acts. There's no way to designate one component "the secure part" and relax about the rest, because the data flows between all of them. The useful image is a secure room: once a confidential file is on the table, everything on the table is confidential until someone deliberately produces a clean, reviewed version to carry out.

A small or narrow-looking output is not automatically safe. A bare "PASS/FAIL" can still disclose that a particular condition exists, betray a party's position, or let an observer infer privileged analysis by watching which inputs produce which verdicts over time. Output shape doesn't determine sensitivity --- what it reveals, including by inference, does.

Standard tenant isolation is the wrong boundary for internal barriers. Most enterprise products isolate one customer's data from another's --- which satisfies standard certifications but treats an entire organization as one trusted space. The barrier that actually matters, for an organization with internal information barriers, runs \emph{inside} the organization. The question isn't "is our data isolated from other customers" but "does the system preserve our internal permission structure everywhere the data and its derivatives travel." Four places are worth checking specifically, because each is a documented way derived data can escape a permission boundary that the front door itself enforced: search indexes and vector databases; embeddings and prompt caches; fine-tuning or training datasets; centralized logs, exports, and telemetry.

These are real threat classes, worth stating with appropriate precision. Researchers have demonstrated reconstruction (inversion) attacks against text embeddings under specific conditions, timing-based side channels related to prompt caching in certain architectures, and gaps in fine-grained access control for some vector-database implementations. These are documented risks with stated prerequisites --- specific architectures, specific attack models --- not consequences that occur automatically in every system. Treat each one as \textbf{a threat model to test for in a specific deployment}, not a certainty to assert about every product.

Privilege loss is a distinct hazard from confidentiality loss --- narrower but graver when it occurs. Privilege and work-product treatment depend on jurisdiction, the role and necessity of the third party, the confidentiality arrangements in place, the purpose of the disclosure, and the precautions taken --- and outcomes have varied materially across recent cases addressing AI-tool use specifically. The defensible statement is only that \textbf{privilege consequences require jurisdiction- and fact-specific legal analysis}, that privilege and work-product protection are distinct legal concepts that shouldn't be collapsed together, and that this is squarely a question for a qualified professional, not an engineering assumption in either direction.

On reconstructability: exact byte-for-byte reproducibility isn't guaranteed by most hosted APIs, but reproducibility improves substantially --- without necessarily becoming exact --- when the model build, runtime, inputs, retrieved-context snapshot, and parameters are frozen and logged; some providers also offer seed-based mechanisms aimed at closer determinism, though typically on a best-effort basis rather than a guarantee. The useful distinction is between \textbf{reconstructing the process} (achievable, and worth requiring as a matter of practice) and \textbf{reproducing the exact output byte-for-byte} (not reliably achievable, and not the actual requirement most governance needs are trying to satisfy).

Two further governance dimensions round this out: \emph{data minimization} is a distinct question from path-tracing --- personal data carries lawful-basis, individual-rights, and cross-border obligations a path-and-boundary analysis doesn't discharge on its own. And \emph{knowledge ownership} --- who ends up owning an organization's own playbooks, taxonomies, and evaluation sets, which can become trapped inside a tooling vendor's proprietary index. The exit question --- can this be exported in an open format on termination --- is worth asking before signing.
